
\section{The problem at a fixed profile, and statement of results}\label{sec:problem}

Sections~\ref{sec:model}--\ref{sec:binomial} set up the model, the problem
at a fixed profile (no step changes the profile) and the algebraic language
it compiles to. Section~\ref{sec:results} states the three results.
% gloss: at a fixed profile

% ---------------------------------------------------------------------------
\subsection{The model}\label{sec:model}

The model is built so that each label acts on each probe by a single
number. Fix nonempty finite sets $M$ of labels, $X$ of base states and $R$
of register values. The probes $\chi_i:X\to\{-1,1\}$, $i\in I$, are
balanced and pairwise orthogonal:
\[
  \sum_{g\in X}\chi_i(g)=0,\qquad
  \sum_{g\in X}\chi_i(g)\chi_{i'}(g)=|X|\,\mathbf 1_{i=i'} .
\]
% gloss: probe
Nontrivial characters of a finite elementary abelian $2$-group are examples,
but only these identities are used. Each label $m$ has a label probability
$\theta_m$, an eigenvalue $\alpha_{mi}$ on each probe and register tables
$L_m(g,h;\cdot,\cdot)$, and
\[
  Q_m(g,h)=\frac{1+\sum_i\alpha_{mi}\chi_i(g)\chi_i(h)}{|X|},\qquad
  P_m((g,r),(h,r'))=\theta_m\,Q_m(g,h)\,L_m(g,h;r,r')
\]
are its base kernel and transition table. Each $Q_m$ is a
\defterm{parity mixture}: it mixes a uniform refresh with moves that keep or
flip probe signs. So every probe is an eigenvector of every base kernel,
with eigenvalue~$\alpha_{mi}$. Here $\alpha_{mi}$ is an eigenvalue of $Q_m$,
not in general of~$P_m$.

Two bounds keep the label probabilities and base kernels positive as
eigenvalues are damped. The \defterm{$\ell^1$ budget} $\kappa<1$ requires
$\sum_i|\alpha_{mi}|\le\kappa$ for every label, while the eigenvalues
themselves may be negative or zero. This budget makes $Q_m$ positive: with
probability $1-\sum_i|\alpha_{mi}|$ the base state refreshes uniformly,
and with probability $|\alpha_{mi}|$ moves to a uniform $h$ with
$\chi_i(h)=\sgn(\alpha_{mi})\chi_i(g)$ (omitting zero-weight terms).
Hence $Q_m(g,h)\ge(1-\kappa)/|X|>0$. The budget survives every decrease in
absolute values, though not every coordinatewise decrease. The probability
floor requires $\theta_m\ge\delta>0$ and
$\sum_m\theta_m=1$; it lets steps in the transition tables bound steps in
the eigenvalues (Section~\ref{sec:tractable}).

The register is hidden from the word tests, which sum it out; it carries
the discrete table-support and level data of the profile
(Section~\ref{sec:profiles}). Its tables $L_m(g,h;\cdot,\cdot)$ are
stochastic, may depend on the base transition $(g,h)$ and keep stationary a
fixed rational register distribution $\nu$ of full support. The register,
levels, decisions and admissible sets
(Sections~\ref{sec:words}--\ref{sec:profiles}) enter the definition of
feasibility and govern which profile steps are possible. All Markov
hardness constructions in this paper use a one-value register, only
level~$0$, and the full word-test list as the sole admissible set.

A \defterm{reversal} is an involution $m\mapsto\bar m$ of the labels,
included so that the stationary time reversal of each parameter vector
remains in the model. It imposes $\alpha_{\bar mi}=\alpha_{mi}$ on every
parameter vector, so these reversal equations must be included in the
eigenvalue constraints $\Phi_p$ (Section~\ref{sec:binomial}). With the
reciprocity $\nu(r)L_m(g,h;r,r')=\nu(r')L_{\bar m}(h,g;r',r)$, the
distribution $\mu(g,r)=\nu(r)/|X|$ is stationary, and
\[
  \mu(g,r)\,P_m((g,r),(h,r'))\,\theta_{\bar m}
  =\mu(h,r')\,P_{\bar m}((h,r'),(g,r))\,\theta_m
\]
even when $\theta_m\neq\theta_{\bar m}$. The stationary time reversal, with
each reversed transition relabeled by the involution, has the same base and
register kernels and label probabilities
$\theta^{\mathrm{rev}}_m=\theta_{\bar m}$. All Markov hardness constructions
use the identity reversal, while the $x^2=\bar xz$ example after
Theorem~\ref{thm:fixed-profile} shows why the reversal equations matter
for~$\Phi_p$.

\begin{definition}[Shared-probe model]\label{def:model}
A \emph{shared-probe model} fixes finite nonempty sets $M$, $X$, $R$,
probes $(\chi_i)_{i\in I}$ as above, a reversal $m\mapsto\bar m$, a rational
distribution $\nu$ on $R$ of full support, a maximum word length $H\ge1$
and rationals $\delta>0$ and $0\le\kappa<1$. Its parameter vectors are the
rational triples $(\theta,\alpha,L)$ with $\theta_m\ge\delta$,
$\sum_m\theta_m=1$, $\sum_i|\alpha_{mi}|\le\kappa$,
$\alpha_{\bar mi}=\alpha_{mi}$, and stochastic tables
$L_m(g,h;\cdot,\cdot)$ that keep $\nu$ stationary and satisfy the
reciprocity above.
\end{definition}

In other words, a parameter vector is a transition-labeled Markov chain
on $X\times R$~\cite[Definition~1]{chenkiefer2014} with parity-mixture
base kernels.
% gloss: labeled Markov chain
In the staged example, the base is
$X=\{0,1\}$ with one probe, $\chi(0)=1$ and $\chi(1)=-1$. Its four labels
$\mathsf x,\mathsf y,\mathsf z,\mathsf f$ each have probability $1/4$, are
their own reversals and have eigenvalues $x,y,z,f$, with $f=0$ along the revision. The label $\mathsf f$ carries the
remaining quarter of the label probability and appears in neither
% gloss: held test
held test. Here $\delta=1/8$, $\kappa=1/2$ and $H=2$. The register has four values
and uniform $\nu$, and every table entry is $1/4$, except that for
$\mathsf z$ at $g=h=0$ the block on values $\{0,1\}$ is
$\operatorname{diag}(1/2,1/2)$.
At label~$\mathsf y$, the conditional probability of keeping the probe
sign changes from $5/8$ to $3/8$ along the revision. Jointly with that
label, it changes from $5/32$ to $3/32$.

% ---------------------------------------------------------------------------
\subsection{Words and tests}\label{sec:words}

A word $w$ is a nonempty string of labels of
length $|w|\le H$, with $|w|_m$ occurrences of~$m$. Its word probability
and word eigenvalues are $\theta(w)=\prod_m\theta_m^{|w|_m}$ and
$\alpha_i(w)=\prod_m\alpha_{mi}^{|w|_m}$. Summing out the register, the
probability of observing $w$ from state $(g,r)$ and ending at base state
$h$ is
\begin{equation}\label{eq:word-joint}
  \theta(w)\,(Q_{w_1}\cdots Q_{w_{|w|}})(g,h)
  =\frac{\theta(w)}{|X|}\Bigl(1+\sum_{i\in I}\alpha_i(w)\chi_i(g)\chi_i(h)\Bigr).
\end{equation}

% gloss: value test
\begin{definition}[Word tests and value tests]\label{def:tests}
A \emph{word test} is $w\sim w'$ (word-against-word) or $w\sim\{w',w''\}$
with $w'\neq w''$ (word-against-pair), all words of one positive length at
most~$H$. A parameter vector passes it if, from every state $(g,r)$ and for
every final base state $h$, the probability~\eqref{eq:word-joint} for $w$
equals that for $w'$, or the sum of those for $w'$ and $w''$. A
\emph{value test} is a named calibration $\alpha_{mi}=c$ with $c\in\QQ$, or
a probability-ratio test $\theta(w)=c\,\theta(w')$ with $|w|=|w'|$ and
$c\in\QQ_{>0}$. A shared-probe model also fixes a final level $n_F\ge0$,
finite lists of word tests and value tests, and the
% gloss: admissible sets of word tests
\emph{admissible sets of word tests}: the subsets of the word-test list
that a configuration may hold together.
\end{definition}

Admissible sets let a model restrict which combinations of word tests may
be held simultaneously. In all Markov hardness constructions, the full
word-test list is the only admissible set.

\begin{lemma}[What a word test says]\label{lem:compile}
For every parameter vector:
(a)~$w\sim w'$ passes if and only if $\theta(w)=\theta(w')$ and
$\alpha_i(w)=\alpha_i(w')$ for every $i\in I$;
(b)~$w\sim\{w',w''\}$ passes if and only if
$\theta(w)=\theta(w')+\theta(w'')$ and
$\theta(w)\alpha_i(w)=\theta(w')\alpha_i(w')+\theta(w'')\alpha_i(w'')$ for
every $i\in I$.
\end{lemma}

Thus a word-against-word test gives one homogeneous rate equation and one
binomial probe equation per probe. A word-against-pair test gives $1+|I|$
displayed equations, possibly dependent. Its probe equations can couple
label probabilities and eigenvalues and need not be binomial. By Lemma~\ref{lem:compile},
the tests of the two examples also give the rate equations
$\theta_{\mathsf x}\theta_{\mathsf y}=\theta_{\mathsf x}\theta_{\mathsf z}$
and $\theta_{\mathsf y}^2=\theta_{\mathsf z}^2$. Lemma~\ref{lem:compile}
and~\eqref{eq:word-joint} are proved in Appendix~\ref{app:problem-proofs}.

% ---------------------------------------------------------------------------
\subsection{Profiles and the fixed-profile problem}\label{sec:profiles}

Reachability is between configurations, whose profile carries the discrete
data.

\begin{definition}[Profiles and feasibility]\label{def:profile}
A \emph{profile} $p=(\Gamma_p,F^p,\mathcal T_p)$ has three parts. The
table support $\Gamma_p$ is a set of indices $(m,g,r,h,r')$. The level
partitions $F^p_0,\dots,F^p_{n_F}$ of $R$ each refine the one before. The
\defterm{held tests} $\mathcal T_p$ are word tests and value tests that
must keep passing. A decision at level $\ell$ is a $0/1$ function on $R$
that is constant on the blocks of $F^p_\ell$. A configuration $A$ is a
profile $p_A$ with a rational parameter vector
$(\theta_A,\alpha_A,L_A)$, equivalently with its transition tables~$P_A$.
It is \emph{feasible} if the positive entries of $P_A$ are exactly
$\Gamma_{p_A}$, every decision $d$ at every level satisfies
\[
  \sum_{r':\,d(r')=1}L_m(g,h;r,r')=\sum_{r':\,d(r')=1}\nu(r')
  \qquad\text{for all }m,g,h,r,
\]
the parameter vector passes every test in $\mathcal T_{p_A}$, and the held
word tests form an admissible set.
\end{definition}

The displayed block-mass condition says that, for every label~$m$, base
move $(g,h)$ and current register value~$r$, the new register value
satisfies each level-$\ell$ decision with the same probability as
under~$\nu$. Because the partitions are nested, any decision
available at an earlier level is also available at level~$n_F$. It
therefore suffices to check the displayed block-mass condition at
level~$n_F$.

% gloss: attained
In other words, the profile records which transitions occur, which register
partitions constrain block masses at each level and which tests are held. A
profile is \defterm{attained} when some feasible rational configuration has
it. In the staged example, levels $0$ and $1$ both group the register
values as $\{0,1\}\mid\{2,3\}$, and the profile holds
$\mathsf{xy}\sim\mathsf{xz}$, admissible on its own, and the calibration
$z=1/4$.

% gloss: fine reachability
\begin{center}
\fbox{\begin{minipage}{0.93\linewidth}
\textsc{Fine Reachability at a fixed profile}\par\smallskip
\emph{Input:} the explicitly presented model, a profile and two rational
parameter vectors at it, in binary.\par\smallskip
\emph{Promise:} both configurations are feasible.\par\smallskip
\emph{Question:} is there, at every positive tolerance, a finite
reconfiguration sequence of feasible rational configurations joining the
two supplied configurations at that profile whose
steps change each transition-table entry $P_m((g,r),(h,r'))$ by less than
the tolerance? Zero steps are allowed when they coincide.\par\smallskip
\emph{Supplied-tolerance variant:} the input also contains
$\varepsilon_0\in\QQ_{>0}$, and the question is whether such a sequence
exists at tolerance~$\varepsilon_0$.
\end{minipage}}
\end{center}

If the answer is yes, the second configuration~$B$ is
\defterm{finely reachable} from the first,~$A$. Real and rational
tolerances give the same verdict.

\begin{remark}[Inputs and equal profiles]\label{rem:inputs}
(i)~Numbers are binary rationals, and the model is given by explicit lists
and tables. The promise covers admissibility and attainment, and no
algorithm here checks it.
(ii)~Profile steps (Section~\ref{sec:directed}) go up a partial order, so a
sequence between configurations at one profile contains no profile step, at
any tolerance. Hence the lower bounds of Theorem~\ref{thm:boundary} hold
with profile steps allowed.
\end{remark}

% ---------------------------------------------------------------------------
\subsection{Binomial systems, peeling and the signed core}\label{sec:binomial}
% gloss: signed core

Held word-against-word tests, calibrations and the fixed reversal equations
give binomial eigenvalue equations, whose coordinates can be zero or
negative; probability-ratio tests give only rate equations.
Section~\ref{sec:core} decides reachability by tracking which coordinates
can vanish together, and with which signs.

% gloss: budgeted signed-binomial system
\begin{definition}[Budgeted signed-binomial systems, peeling, signed core]\label{def:system}
Let $J$ be finite.
\begin{enumerate}
\item[(a)] A \emph{binomial test} is $x^a=c\,x^b$ with $a,b\in\NN^J$ and
$c\in\QQ$, where $x^a=\prod_jx_j^{a_j}$ and $x_j^0=1$ also at $x_j=0$. An
\emph{$\ell^1$ budget} is $\sum_{j\in J_\beta}|x_j|\le\kappa_\beta$ for an
explicitly listed subset $J_\beta\subseteq J$ and a rational $\kappa_\beta$.
A \defterm{budgeted signed-binomial system} is a finite test list $\Phi$
with finitely many budgets; its feasible points are those of $\RR^J$ that
satisfy them all.
\item[(b)] The \emph{signed support} of $x$ is $\supp x$ with the sign of each
$x_j$ there. A test is \emph{active} on $S\subseteq J$ when
$\supp a\cup\supp b\subseteq S$, and $\Delta_S$ has one row $a-b$,
restricted to the columns in~$S$, for each test active on~$S$.
\item[(c)] A \defterm{weight vector} $\eta\in\NN^S$ with $\Delta_S\eta=0$ defines
the \emph{peeling step} $S\to T$ with $T=\{j\in S:\eta_j=0\}$. For a point
$x$ of support $S$, its \defterm{monomial curve} is $x_jt^{\eta_j}$ on $S$
and zero off $S$, with $t^0=1$, as $t\downarrow0$. $S$ \emph{peels to} $T$
if zero or more peeling steps lead from $S$ to~$T$, with the active tests
recomputed at each support.
\item[(d)] The \emph{core} $\core_\Phi(S)$ is the least support that $S$ peels to
(Lemma~\ref{lem:restriction}), and the \defterm{signed core}
$\SC_\Phi(x)$ equals $\sgn x_j$ for $j\in\core_\Phi(\supp x)$ and $0$
elsewhere ($\Phi$ is omitted when fixed).
\item[(e)] Let $p$ be a profile of a model whose word tests are all
word-against-word. Its \emph{eigenvalue constraints} $\Phi_p$ on
$J=M\times I$ are one reversal equation $\alpha_{\bar mi}=\alpha_{mi}$ per
coordinate, one equation $\alpha_i(w)=\alpha_i(w')$ per held
word-against-word test and probe, and one equation $\alpha_{mi}=c$ per held
calibration. Its budgets are $\sum_i|\alpha_{mi}|\le\kappa$, one per label.
We write $\SC_p$ for $\SC_{\Phi_p}$.
\end{enumerate}
\end{definition}

In other words, positive weights send their coordinates to zero while the
active tests stay satisfied, and the signed core records the surviving
support and signs, not magnitudes. The all-zero weight vector gives the
identity step. On the signed support of a feasible point, a test with
$c\neq0$ is inactive exactly when both its monomials vanish; for $c=0$ the
left monomial must vanish, so the test is inactive. In $\Phi_p$ the case
$c=0$ occurs only for a zero calibration. In the staged example, $\Phi_p$ is
$xy=xz$, $z=1/4$ and identity reversal equations, and both endpoints have
signed core $z$ alone, positive (Section~\ref{sec:restriction}).

Write $\FQ$ for the set of feasible rational points of a budgeted
signed-binomial system.

\begin{center}
\fbox{\begin{minipage}{0.93\linewidth}
\textsc{Fine Connectivity}\par\smallskip
\emph{Input:} a budgeted signed-binomial system (coordinates $J$, tests
$x^a=c\,x^b$ with nonnegative integer exponents encoded in binary, with
sparse lists allowed, and rational $c$,
budgets $\sum_{j\in J_\beta}|x_j|\le\kappa_\beta$), and two rational
points $x,x'$.\par\smallskip
\emph{Promise:} $\FQ$ is bounded, and $x$ and $x'$ are feasible.\par\smallskip
\emph{Question:} for every rational $\varepsilon>0$, is there a finite
sequence of feasible rational points from $x$ to $x'$ in which consecutive
points differ by less than $\varepsilon$ in every coordinate?\par\smallskip
% gloss: affine trinomial
\emph{With affine trinomials:} the input may also list affine trinomials
$x_j+y_j=k$. The normalized version, with equations $x_{j0}+x_{j1}=1$, is
the restriction on which hardness already holds.
\end{minipage}}
\end{center}

% ---------------------------------------------------------------------------
\subsection{Statement of results}\label{sec:results}

% gloss: P-complete
\begin{theorem}[The boundary, at a fixed profile]\label{thm:boundary}
Consider Fine Reachability at a fixed profile for explicitly presented
shared-probe models, with binary rational data and two configurations
promised feasible.
\begin{enumerate}
\item[(i)] If every word test compares one word with one word, the problem
is P-complete under deterministic log-space many-one reductions, with
polynomial-size, polynomial-time checkable certificates for both answers.
% gloss: PSPACE-hard
\item[(ii)] If exactly one word test compares a single word with the sum
over two distinct words, and every other word test compares one word with
one word, the problem is PSPACE-hard under polynomial-time many-one
reductions. The same reduction proves PSPACE-hardness of sequence existence
at a supplied rational tolerance $\varepsilon_0>0$ whose reciprocal is
polynomially bounded in the length of the reduction's input.
\end{enumerate}
Both lower bounds hold on instances with one common profile, at which every
test is held, the table support is full and only level~$0$ occurs, and with
a one-value register, identity reversal and $\kappa=1/2$. The instances for~(i) have a binary base, one probe,
one calibration and every word in a word test of length three. Those
for~(ii) have the word-against-pair test of length one or two, every other
word test of length three, and distinct endpoints.
\end{theorem}

% gloss: exactly one word-against-pair test
Thus the single-word problem is P-complete, and PSPACE-hardness already
holds with exactly one word-against-pair test. We prove no upper bound for the full problem of part~(ii):
a PSPACE upper bound would put the existence of rational zeros of integer
polynomials in PSPACE (Theorem~\ref{thm:rational-roots}). On each output
of the reduction of Section~\ref{sec:hard} from nondeterministic constraint
logic (NCL)~\cite{hearndemaine}, the NCL verdict agrees with sequence
existence at~$\varepsilon_0$ and at every smaller positive tolerance. This
equivalence can fail at larger tolerances. Indeed, at tolerance~$2$, any
two feasible configurations at one profile are one step apart, yet an NCL
graph with two AND vertices joined by parallel edges of weights $2$, $1$
and~$1$ has two legal orientations, each isolated.

% gloss: G9
\begin{theorem}[Fixed-profile criterion]\label{thm:fixed-profile}
Fix a shared-probe model in which every word test is word-against-word, and
let $A$ and $B$ be feasible rational configurations at a common
profile~$p$. Then $B$ is finely reachable from $A$ if and only if
$\SC_p(\alpha_A)=\SC_p(\alpha_B)$, where $\SC_p$ is the signed core for the
eigenvalue constraints $\Phi_p$, including the reversal equations.
\end{theorem}

Thus, in the word-against-word class at a fixed profile, fine
reachability depends only on the support that survives peeling and the
signs there. In particular, the criterion joins the staged
example's endpoints, whose signed cores agree, and separates the sign
example's, whose cores differ at~$y$. The reversal equations cannot be
omitted. Take a two-state base with one probe, a one-value register,
$\kappa=1/2$ and labels $\mathsf x,\bar{\mathsf x},\mathsf z$ of probability
$1/3$, where $\bar{\mathsf x}$ reverses~$\mathsf x$. Hold
$\mathsf{xx}\sim\bar{\mathsf x}\mathsf z$ and $z=1/4$, so that $x^2=\bar xz$.
Without the reversal equation, the weights $(1,2,0)$ peel the feasible
endpoint $(x,\bar x,z)=(1/4,1/4,1/4)$ to the support of the feasible
endpoint $(0,0,1/4)$, so the computed cores agree. With it, neither $x$
nor $\bar x$ can peel, so the cores differ. Theorem~\ref{thm:directed}
extends the criterion to profile steps: it adds $p_A\preceq p_B$ and uses
the source profile's constraints.

\begin{corollary}[The algebraic form]\label{cor:algebraic}
Consider explicitly binary-encoded budgeted signed-binomial systems whose
feasible rational set is promised bounded, with two supplied feasible
rational points. Deciding whether, for every rational $\varepsilon>0$, a
finite sequence of feasible rational points joins them with every
coordinate changing by less than $\varepsilon$ per step is P-complete under
deterministic log-space many-one reductions whose outputs meet the
promises. Allowing a growing number of scalar affine trinomials
$x_j+y_j=k$ makes it PSPACE-hard under polynomial-time many-one reductions
with the same promises, even with $k=1$ throughout.
\end{corollary}

% gloss: affine trinomial
Thus budgeted signed-binomial Fine Connectivity
(Section~\ref{sec:binomial}) is P-complete, and allowing a
growing number of scalar affine trinomials gives PSPACE-hardness under the
same promises. %
The signed-core characterization used for the upper bound requires
rational boundedness. Without it, points with different signed cores can
be finely connected, that is, joined at every tolerance, as the example in
Section~\ref{sec:classification} shows.
% gloss: finely connected
For this rational problem the paper gives no upper bound with trinomials,
and hardness with one trinomial is open.
